% ---------- ABSTRACT (spec section 4) : page iii ----------
% Rewritten 2026-09-05 to CSE_Thesis_Final_Abstract_Guidelines.pdf:
%   * 200-300 words (was 414)
%   * flow Background/Problem -> Objective -> Methodology -> Key Results ->
%     Contribution/Conclusion
%   * only the essential numbers (7, was about 21) -- the supervisor's note was
%     that too many results were reported
%   * past tense for methods and completed experiments, present tense for the
%     conclusion
%   * no citations, equations, or implementation detail; no limitations or
%     future work
% Revised again for cohesion: the findings were previously four free-standing
% sentences. Each now carries the argument to the next -- the second phase is
% given its reason, the architectural results are framed as one set, the
% assignment result is set against them as a contrast, and the real-imagery
% result follows from it.
\phantomsection\addcontentsline{toc}{fmatter}{Abstract}
\frontheading{Abstract}

Finding people in aerial disaster imagery is a search-and-rescue problem in
which the targets are often only a few pixels wide, partly occluded, and set
against rubble, water, or smoke. Detectors tuned for ordinary object sizes lose
them, and published work on the task reports strong aggregate scores without
isolating which design change earns them. This thesis presents a compact
real-time detector for that regime together with the evidence for each of its
parts: four configurations were trained under a single frozen protocol on the
C2A disaster-imagery benchmark, beginning from a YOLO11m baseline and adding,
one change at a time, a channel-and-spatial attention module, a stride-4
detection scale, and a bidirectional state-space neck. Because a benchmark
result can be an artefact of how the data is divided, those conclusions were
then retested on a scene-disjoint partition of the same images, under a
modified label-assignment rule, and on real footage from drone flights, campus
scenes, and disaster news video. Of the three architectural changes only the
stride-4 scale mattered, raising $\mathrm{AP}_{50}$ from $0.843$ to $0.853$;
the state-space neck cost a factor of $2.8$ in latency for no accuracy gain and
was rejected. The larger improvement on the smallest targets came not from
architecture but from supervision, as blending a scale-tolerant similarity into
label assignment raised recall below eight pixels by $1.6$ points across
replicated seeds while leaving the coarser size bands unchanged. That advantage
did not carry to real imagery unaided, where the synthetically trained model
reached $0.335$~$\mathrm{AP}_{50}$ until adaptation on a few hundred labelled
frames lifted it to $0.795$. The resulting detector uses $19.6$~million
parameters and runs in $14.6$~ms per image, and the wider lesson is that at
this target size the supervision signal is worth more than added network
capacity.
\clearpage
